
This paper investigates the effect of error feedback granularity on LLM code repair success. In a controlled experiment using CodeLlama-7B-Instruct on 304 MBPP runtime error cases, five granularity levels were compared: pass/fail only (G0), error type (G1), error message (G2), error with line number (G3), and full stack trace (G4). One-way ANOVA indicates that granularity has a statistically significant effect on repair success (F=23.89, p=3.5e-19, eta-squared=0.059). Contrary to the initial hypothesis that intermediate granularity would be optimal, the results show that minimal feedback (G0, G1) achieves higher repair success rates (41.8\% and 40.8\%) than detailed feedback (G2: 18.4\%, G3: 16.8\%, G4: 22.7\%). The difference between G0 and G3 is 25 percentage points in favor of G0 (McNemar p=5.2e-22). These findings are specific to CodeLlama-7B-Instruct on MBPP and should not be extrapolated to larger models or different benchmarks without further empirical validation.
